% Options for packages loaded elsewhere
\PassOptionsToPackage{unicode}{hyperref}
\PassOptionsToPackage{hyphens}{url}
\PassOptionsToPackage{dvipsnames,svgnames,x11names}{xcolor}
%
\documentclass[
  letterpaper,
]{article}

\usepackage{amsmath,amssymb}
\usepackage{iftex}
\ifPDFTeX
  \usepackage[T1]{fontenc}
  \usepackage[utf8]{inputenc}
  \usepackage{textcomp} % provide euro and other symbols
\else % if luatex or xetex
  \usepackage{unicode-math}
  \defaultfontfeatures{Scale=MatchLowercase}
  \defaultfontfeatures[\rmfamily]{Ligatures=TeX,Scale=1}
\fi
\usepackage{lmodern}
\ifPDFTeX\else  
    % xetex/luatex font selection
\fi
% Use upquote if available, for straight quotes in verbatim environments
\IfFileExists{upquote.sty}{\usepackage{upquote}}{}
\IfFileExists{microtype.sty}{% use microtype if available
  \usepackage[]{microtype}
  \UseMicrotypeSet[protrusion]{basicmath} % disable protrusion for tt fonts
}{}
\makeatletter
\@ifundefined{KOMAClassName}{% if non-KOMA class
  \IfFileExists{parskip.sty}{%
    \usepackage{parskip}
  }{% else
    \setlength{\parindent}{0pt}
    \setlength{\parskip}{6pt plus 2pt minus 1pt}}
}{% if KOMA class
  \KOMAoptions{parskip=half}}
\makeatother
\usepackage{xcolor}
\usepackage[margin=1in]{geometry}
\setlength{\emergencystretch}{3em} % prevent overfull lines
\setcounter{secnumdepth}{5}
% Make \paragraph and \subparagraph free-standing
\makeatletter
\ifx\paragraph\undefined\else
  \let\oldparagraph\paragraph
  \renewcommand{\paragraph}{
    \@ifstar
      \xxxParagraphStar
      \xxxParagraphNoStar
  }
  \newcommand{\xxxParagraphStar}[1]{\oldparagraph*{#1}\mbox{}}
  \newcommand{\xxxParagraphNoStar}[1]{\oldparagraph{#1}\mbox{}}
\fi
\ifx\subparagraph\undefined\else
  \let\oldsubparagraph\subparagraph
  \renewcommand{\subparagraph}{
    \@ifstar
      \xxxSubParagraphStar
      \xxxSubParagraphNoStar
  }
  \newcommand{\xxxSubParagraphStar}[1]{\oldsubparagraph*{#1}\mbox{}}
  \newcommand{\xxxSubParagraphNoStar}[1]{\oldsubparagraph{#1}\mbox{}}
\fi
\makeatother


\providecommand{\tightlist}{%
  \setlength{\itemsep}{0pt}\setlength{\parskip}{0pt}}\usepackage{longtable,booktabs,array}
\usepackage{calc} % for calculating minipage widths
% Correct order of tables after \paragraph or \subparagraph
\usepackage{etoolbox}
\makeatletter
\patchcmd\longtable{\par}{\if@noskipsec\mbox{}\fi\par}{}{}
\makeatother
% Allow footnotes in longtable head/foot
\IfFileExists{footnotehyper.sty}{\usepackage{footnotehyper}}{\usepackage{footnote}}
\makesavenoteenv{longtable}
\usepackage{graphicx}
\makeatletter
\newsavebox\pandoc@box
\newcommand*\pandocbounded[1]{% scales image to fit in text height/width
  \sbox\pandoc@box{#1}%
  \Gscale@div\@tempa{\textheight}{\dimexpr\ht\pandoc@box+\dp\pandoc@box\relax}%
  \Gscale@div\@tempb{\linewidth}{\wd\pandoc@box}%
  \ifdim\@tempb\p@<\@tempa\p@\let\@tempa\@tempb\fi% select the smaller of both
  \ifdim\@tempa\p@<\p@\scalebox{\@tempa}{\usebox\pandoc@box}%
  \else\usebox{\pandoc@box}%
  \fi%
}
% Set default figure placement to htbp
\def\fps@figure{htbp}
\makeatother
% definitions for citeproc citations
\NewDocumentCommand\citeproctext{}{}
\NewDocumentCommand\citeproc{mm}{%
  \begingroup\def\citeproctext{#2}\cite{#1}\endgroup}
\makeatletter
 % allow citations to break across lines
 \let\@cite@ofmt\@firstofone
 % avoid brackets around text for \cite:
 \def\@biblabel#1{}
 \def\@cite#1#2{{#1\if@tempswa , #2\fi}}
\makeatother
\newlength{\cslhangindent}
\setlength{\cslhangindent}{1.5em}
\newlength{\csllabelwidth}
\setlength{\csllabelwidth}{3em}
\newenvironment{CSLReferences}[2] % #1 hanging-indent, #2 entry-spacing
 {\begin{list}{}{%
  \setlength{\itemindent}{0pt}
  \setlength{\leftmargin}{0pt}
  \setlength{\parsep}{0pt}
  % turn on hanging indent if param 1 is 1
  \ifodd #1
   \setlength{\leftmargin}{\cslhangindent}
   \setlength{\itemindent}{-1\cslhangindent}
  \fi
  % set entry spacing
  \setlength{\itemsep}{#2\baselineskip}}}
 {\end{list}}
\usepackage{calc}
\newcommand{\CSLBlock}[1]{\hfill\break\parbox[t]{\linewidth}{\strut\ignorespaces#1\strut}}
\newcommand{\CSLLeftMargin}[1]{\parbox[t]{\csllabelwidth}{\strut#1\strut}}
\newcommand{\CSLRightInline}[1]{\parbox[t]{\linewidth - \csllabelwidth}{\strut#1\strut}}
\newcommand{\CSLIndent}[1]{\hspace{\cslhangindent}#1}

\usepackage{booktabs}
\usepackage{caption}
\usepackage{longtable}
\usepackage{colortbl}
\usepackage{array}
\usepackage{anyfontsize}
\usepackage{multirow}
\usepackage{setspace}
\doublespacing
\usepackage{lineno}
\linenumbers
\usepackage{authblk}
\usepackage{marvosym}
\author[1]{Dominic Woolf\,\thanks{Correspondence to d.woolf@cornell.edu}}
\author[2]{Jocelyn Lavallee}
\author[2]{Anna Raffeld}
\author[2]{Doria Gordon}
\author[2]{Emily Oldfield}
\affil[1]{Cornell University}
\affil[2]{Environmental Defense Fund}
\renewcommand\author[2][]{}
\makeatletter
\@ifpackageloaded{caption}{}{\usepackage{caption}}
\AtBeginDocument{%
\ifdefined\contentsname
  \renewcommand*\contentsname{Table of contents}
\else
  \newcommand\contentsname{Table of contents}
\fi
\ifdefined\listfigurename
  \renewcommand*\listfigurename{List of Figures}
\else
  \newcommand\listfigurename{List of Figures}
\fi
\ifdefined\listtablename
  \renewcommand*\listtablename{List of Tables}
\else
  \newcommand\listtablename{List of Tables}
\fi
\ifdefined\figurename
  \renewcommand*\figurename{Figure}
\else
  \newcommand\figurename{Figure}
\fi
\ifdefined\tablename
  \renewcommand*\tablename{Table}
\else
  \newcommand\tablename{Table}
\fi
}
\@ifpackageloaded{float}{}{\usepackage{float}}
\floatstyle{ruled}
\@ifundefined{c@chapter}{\newfloat{codelisting}{h}{lop}}{\newfloat{codelisting}{h}{lop}[chapter]}
\floatname{codelisting}{Listing}
\newcommand*\listoflistings{\listof{codelisting}{List of Listings}}
\makeatother
\makeatletter
\makeatother
\makeatletter
\@ifpackageloaded{caption}{}{\usepackage{caption}}
\@ifpackageloaded{subcaption}{}{\usepackage{subcaption}}
\makeatother

\usepackage{bookmark}

\IfFileExists{xurl.sty}{\usepackage{xurl}}{} % add URL line breaks if available
\urlstyle{same} % disable monospaced font for URLs
\hypersetup{
  pdftitle={Optimal biomass-based climate change mitigation pathways are spatially heterogeneous due to localized technoeconomic and biophysical constraints},
  colorlinks=true,
  linkcolor={blue},
  filecolor={Maroon},
  citecolor={Blue},
  urlcolor={Blue},
  pdfcreator={LaTeX via pandoc}}


\title{Optimal biomass-based climate change mitigation pathways are
spatially heterogeneous due to localized technoeconomic and biophysical
constraints}
\author{}
\date{2026-08-27}

\begin{document}
\maketitle
\begin{abstract}
The use of biomass for climate change mitigation offers substantial
technical potential, but optimal deployment is highly dependent on
spatial, technoeconomic, and biophysical constraints. Blanket climate
policies frequently fail to capture this heterogeneity. Here we apply a
spatially explicit techno-economic and consequential life cycle
assessment framework deployed at a 0.1-degree resolution across the
United States, Europe, India, and China. We evaluate the economic
viability and carbon abatement potential of three distinct biomass
utilization pathways: Bioenergy (BES), Bioenergy with Carbon Capture and
Storage (BECCS), and Bioenergy with Biochar Carbon Storage (BEBCS).
Localized constraints, such as distance to geological
CO\textsubscript{2} sinks, baseline grid carbon intensity, and native
soil cation exchange capacity, dictate the economically optimal pathway.
Traditional BES tends to dominate at low carbon prices, particularly in
regions with high wholesale energy prices like Europe, or high grid
carbon intensities like India. Biochar (BEBCS) emerges as a highly
competitive transition technology at moderate carbon prices, and serves
as the optimal long-term pathway in specific agro-ecological zones with
degraded soils, such as southern China. Conversely, high capital costs
and piecewise long-distance pipeline penalties act as severe barriers to
BECCS, which only achieves universal dominance at high carbon prices
(e.g., \$150/tCO\textsubscript{2}e) where carbon capture revenues
finally overwhelm infrastructure penalties. Ultimately, effective
climate mitigation requires localized policies that match biomass
resources to their most efficient regional end-use.
\end{abstract}


\section{Introduction}\label{introduction}

Biochar has been widely promoted for climate change mitigation, with
co-benefits for soil health and crop production when used as a soil
amendment.(\citeproc{ref-woolf2010}{\emph{1}},
\citeproc{ref-lehmann2021biochar}{\emph{2}}) Globally, biochar has the
sustainable technical potential for carbon sequestration and emissions
reductions of up to 3.4--6.3 Pg CO\textsubscript{2}e. However, the
costs, benefits and tradeoffs for biochar vary spatially, creating
substantial heterogeneity in its feasibility at a given carbon price.
This variability is further affected by the fact that biochar must
compete with alternative uses for the limited supply of available waste
biomass. Indeed, alternative uses of biomass, such as producing
bioenergy with or without carbon capture and sequestration can also help
to mitigate climate change and may provide greater climate mitigation,
lower costs, and/or greater co-benefits in certain
contexts.(\citeproc{ref-woolf2016}{\emph{3}}) This depends on a number
of different factors, including soil type, cropland productivity,
climate, energy fuel mix, feedstock availability, distance to
CO\textsubscript{2} sinks for carbon capture and storage, energy prices
and food prices. For example, bioenergy production may provide a greater
climate benefit in regions that are highly dependent on carbon-intensive
energy sources and have high soil fertility, whereas biochar is more
favorable where low-fertility soils can benefit from its soil-health
improvement.(\citeproc{ref-woolf2016}{\emph{3}})

It is, therefore, critical to understand where biochar could be the best
use of resources from a climate mitigation perspective. Currently, most
biochar assessments focus on the technical potential, without mapping
out the tradeoffs or identifying optimal pathways based on location.
While there are some biochar life-cycle assessments that account for
alternative biomass uses, they are highly context-specific and difficult
to scale up on(\citeproc{ref-peters2015}{\emph{4}},
\citeproc{ref-lefebvre2021}{\emph{5}}). As such, this project aims to
provide a detailed feasibility and tradeoff analysis across four regions
(North America, India, China, and Europe) to map where biochar
represents the most efficient use of biomass to reduce atmospheric
greenhouse gasses (GHGs). These regions represent a diversity of
geographic, socio-economic, and agro-ecological conditions for a broad
assessment of how constraints, tradeoffs and benefits vary across
regions. This work will allow policy-makers and practitioners to make
scientifically-informed decisions about which mitigation pathways are
best suited for their geographies. It can also help contextualize
biochar's mitigation potential by identifying areas that maximize its
climate benefits and realistically support its production and
application.

\section{Methods}\label{methods}

\subsection{Model Architecture}\label{model-architecture}

The C-SCAPE (Carbon Spatial Cost and Allocation Pathways Evaluator)
model (version 1.0) evaluates three biomass utilization pathways:
Bioenergy (BE), Bioenergy with Carbon Capture and Storage (BECCS), and
Bioenergy with Biochar Carbon Storage (BEBCS). The model evaluates a
spatially-varying net present value (NPV) for each technology using a
discounted cost-benefit analysis, with GHG emissions and removals
converted to 100-year CO\textsubscript{2} equivalent values and
internalized into the economic analysis using an assumed carbon price.

The model was deployed across the coterminous United States, Europe,
India, and China. All spatial inputs were projected into
region-specific, equal-area coordinate reference systems (EPSG:5070,
EPSG:3035, EPSG:8859, and ESRI:102025, for the US, Europe, India, and
China, respectively), and harmonized to a 0.1-degree resolution. The
model was implemented in the R programming language, using the Terra
package for spatial analysis.

The following model description outlines the assumptions,m equations and
parameters used in the model. The default parameter values are provided
in a table at the end of the Methods.

\subsubsection{Biomass Feedstock Sourcing and
Logistics}\label{biomass-feedstock-sourcing-and-logistics}

\paragraph{Biomass Transport Costs and
Emissions}\label{biomass-transport-costs-and-emissions}

High resolution rasters of available crop residue biomass were derived
from Karan et al.(\citeproc{ref-karan2023}{\emph{6}}). The average
transport distance for biomass was calculated from this using focal
statistics based on the annual biomass requirement for the conversion
facility. The required annual biomass mass (\(M_{req}\)) is determined
by the target thermal capacity, the capacity factor, and the biomass
lower heating value (LHV). A spatial focal sum algorithm was applied to
aggregate the surrounding biomass over a discrete series of expanding
radii (from 5 km up to 500 km). The required collection radius was then
calculated via quadratic interpolation between these radial bins. All
domains were uniformly projected and resampled to the 0.1 degree
(\(\sim 10\text{ km}\)) equal-area grid prior to executing the focal sum
algorithms to ensure methodological consistency. Assuming a circular
catchment geometry, the average straight-line collection distance is
approximated as two-thirds of the interpolated bounding radius
(i.e.~assuming homogeneous distribution within the catchment area). The
average straight-line distance was multiplied by a tortuosity factor of
1.3 to calculate the effective road transport
distance.(\citeproc{ref-sultana2014}{\emph{7}}) Biomass logistical costs
were modeled using a fixed material handling and loading cost and a
variable trucking cost applied to the effective road distance. Scope
three transportation emissions associated with feedstock procurement
were estimated using a heavy-duty diesel transport emissions factor
applied to the effective distance and deducted from the net carbon
abatement.

\paragraph{Fuel Quality}\label{fuel-quality}

Agricultural crop residues typically have a higher ash content than the
wood chips more commonly used in bioenergy combustion. High-ash
feedstocks carry significant fouling and slagging risks, requiring
additional capital expenditures (CAPEX) and operating expenditure
(OPEX), for example through deployment of Circulating Fluidized Bed
(CFB) boilers. This fuel quality penalty is modelled as a proportional
(25\%) increase in base CAPEX for combustion technologies (BE and
BECCS), a 50\% increase in operations and maintenance (O\&M) costs, and
a 10\% reduction in overall electrical conversion efficiency, relative
to a wood-chip fired bioenergy
plant.(\citeproc{ref-luan2025a}{\emph{8}}--\citeproc{ref-garcuxeda2017}{\emph{10}})

\paragraph{Regional Pricing
Structures}\label{regional-pricing-structures}

Feedstock pricing is not treated as a static global constant; rather, it
is derived via a spatially explicit shadow-pricing logic converted to a
normalized \(\text{USD/Mg}\) equivalent.

\begin{itemize}
\item
  \textbf{United States:} Establishes a baseline farm-gate cost
  (\(\$70.00/\text{Mg}\)) combined with a nutrient replacement penalty
  (e.g., \(\$25.06/\text{Mg}\) for corn stover) to compensate for the
  export of field nutrients.
\item
  \textbf{Europe:} Utilizes a baseline NUTS-3 regional roadside cost,
  modified by a temporal volatility multiplier that accounts for interim
  storage costs (up to a 36.4\% premium for 6-month storage).
\item
  \textbf{India:} Capitalizes on the necessity to avoid crop stubble
  burning. The model assumes a \(\$0\) stumpage fee, but imposes a
  step-cost for densification: raw bale transport costs apply under 50
  km, while more expensive pelletization costs are triggered for
  catchment zones exceeding 50 km.
\item
  \textbf{China:} Modifies a baseline plant-gate cost with exogenous
  spatial risk multipliers, applying premiums for localized weather
  volatility and rural expansion risks.
\end{itemize}

\subsubsection{Technology Pathways and
Thermodynamics}\label{technology-pathways-and-thermodynamics}

The core thermodynamic engine evaluates three distinct biomass
utilization pathways, processing the standardized
\(250\text{ MW}_{\text{th}}\) biomass input into varying ratios of
exported electrical energy and sequestered carbon.

\paragraph{Bioenergy (BES)}\label{bioenergy-bes}

The BES pathway represents a traditional, dedicated biomass combustion
facility. In this baseline scenario, the feedstock is fully oxidized to
maximize electricity exported to the regional grid. Because there is no
active carbon capture, the physical carbon embedded in the biomass is
returned to the atmosphere.

The technoeconomic parameters for this modern combustion archetype
assume an electrical conversion efficiency of 30\%. Capital expenditures
(CAPEX) are standardized at \(\$3,000/\text{kW}_e\), scaling
non-linearly with plant capacity via a standard 0.7 scaling factor.
Operations and maintenance (O\&M) are assessed as a fixed annual rate of
4\% of total CAPEX. Total electrical output is derived directly from the
biomass lower heating value (LHV, assumed as \(18.6\text{ GJ/Mg}\) for
dry, ash-free feedstock) and the plant's operational capacity factor
(85\%).

\paragraph{Bioenergy with Carbon Capture and Storage
(BECCS)}\label{bioenergy-with-carbon-capture-and-storage-beccs}

The BECCS pathway augments the standard combustion baseline with a
post-combustion carbon capture unit. The addition of this unit imposes a
significant parasitic energy penalty, primarily driven by solvent
regeneration and CO\textsubscript{2} compression requirements.
Consequently, the net electrical conversion efficiency drops from the
30\% BES baseline to 28\%.

The capture process is parameterized with a 90\% capture efficiency
rate. The total volume of captured CO\textsubscript{2} routed to
geological storage is calculated stoichiometrically based on the carbon
fraction of the incoming biomass feed (default 48\% C), converted to
CO\textsubscript{2} mass equivalents using the 44/12 molar ratio. The
increased infrastructure complexity elevates the base CAPEX to
\(\$4,000/\text{kW}_e\), with annual O\&M scaling concurrently to 5\% of
CAPEX.

\paragraph{Bioenergy with Biochar Storage (BEBCS) and Pyrolysis
Physics}\label{bioenergy-with-biochar-storage-bebcs-and-pyrolysis-physics}

Unlike the fully oxidative pathways, the BEBCS pathway relies on
thermochemical pyrolysis, producing a solid, recalcitrant carbon matrix
(biochar) alongside volatile syngas and bio-oil. The model employs a
rigorous mass and energy balance approach constrained by elemental
conservation of Carbon, Hydrogen, and Oxygen (C, H, O).

The mass yield and elemental composition of the resulting biochar are
determined dynamically as a function of the peak pyrolysis temperature
(\(T_{py}\)) agarnd the lignin fraction of the feedstock. Specific
yields for non-condensable gases---hydrogen (\(H_2\)), carbon monoxide
(\(CO\)), methane (\(CH_4\)), and ethylene (\(C_2H_4\))---are calculated
explicitly based on temperature-dependent reaction kinetics.

To close the mass balance, the unaccounted elemental fractions
(\(U_c, U_h, U_o\)) are stoichiometrically partitioned into bio-oil,
water (\(H_2O\)), and carbon dioxide (CO\textsubscript{2}). The model
restricts bio-oil to a fixed empirical composition
(\(62.5\%\text{ C}, 7.56\%\text{ H}, 29.94\%\text{ O}\)) to solve the
final product distribution analytically.

The net usable energy flux for the BEBCS pathway is calculated by
subtracting the process heat requirements from the total enthalpy of the
gaseous and liquid products. The higher heating value (HHV) of the
biochar is calculated via the Dulong formula, utilizing the
temperature-dependent C, H, and O fractions. Heat for the endothermic
pyrolysis reaction is supplied entirely by the combustion of the
volatile syngas and bio-oil fractions; any surplus thermal energy is
converted to electricity using the baseline BES efficiency parameters.
If the syngas and bio-oil are insufficient to meet the thermal load, a
portion of the biochar is oxidized to close the energy balance,
proportionally reducing the final carbon sequestration yield.

\subsubsection{Carbon Transport and Geological
Storage}\label{carbon-transport-and-geological-storage}

\paragraph{Spatial Routing and Terrain
Friction}\label{spatial-routing-and-terrain-friction}

The spatial routing of captured CO\textsubscript{2} is governed by an
economic optimization algorithm that dynamically selects the lowest-cost
transport modality---pipeline or maritime shipping---for every grid
cell. Standard spatial optimization models frequently rely on simple
Euclidean distances or Voronoi tessellations, which artificially
generate circular geographical discontinuities in net present value
(NPV) maps. To eliminate these artifacts and accurately model the
economic constraints of CO\textsubscript{2} transport, this framework
implements a deterministic Lowest Cost Routing logic that relies on a
terrain-optimized Least Cost Path (LCP) algorithm rather than linear
assumptions.

Topographic friction surfaces were generated by deriving the slope (in
degrees) from the Global Digital Elevation Model (GDEM). To preserve
localized linear barriers---such as mountain ridges or steep
ravines---during spatial aggregation to the base model resolution, the
high-resolution slope data was aggregated using a 95th percentile
percolation threshold. The terrain friction surface was then
parameterized using an exponential cost function:

\[M(\theta) = \exp(0.25 \cdot \theta)\]

where \(\theta\) represents the aggregated slope. Furthermore, to ensure
that proposed CO\textsubscript{2} infrastructure avoids ecologically
sensitive or legally restricted regions, protected areas from the World
Database on Protected Areas (WDPA) were rasterized and treated as
absolute barriers. This forces the routing algorithm to find the most
efficient path around protected zones. Once the minimal effective
transport route is mapped, the distances are adjusted by a standard
tortuosity multiplier of 1.25 to account for right-of-way routing and
localized physical avoidance.

\paragraph{Maritime Shipping Modality}\label{maritime-shipping-modality}

For offshore geological sinks, or for highly isolated onshore sources,
CO\textsubscript{2} is transported via maritime or rail networks. This
modality is modeled with a step-cost function dominated by the
capital-intensive fixed entry costs of liquefaction and terminal
buffering, paired with highly efficient variable voyage costs:

\[C_{shipping} = C_{liquefaction} + C_{terminal} + (c_{voyage} \cdot d_{effective})\]

\paragraph{Hub-and-Spoke Pipeline Modality (Piecewise
Linear)}\label{hub-and-spoke-pipeline-modality-piecewise-linear}

Standard length-invariant TEA cost functions for CO\textsubscript{2}
pipelines adequately model regional networks but fundamentally
underestimate the capital and energetic burdens of long-distance
transport. Specifically, maintaining CO\textsubscript{2} in a dense or
supercritical phase over extended distances requires the deployment of
intermediate booster compression stations to overcome frictional
pressure drops.

To simulate this without computationally intensive grid-level pump
topology, pipeline CAPEX is calculated using a piecewise linear
hub-and-spoke framework:

\begin{enumerate}
\def\labelenumi{\arabic{enumi}.}
\item
  \textbf{The Feeder Leg:} Sources build a dedicated feeder pipeline for
  the first 50 km (\(d_{feeder}\)), bearing the full capacity-scaled
  CAPEX.
\item
  \textbf{The Trunk Leg (Base Phase):} For distances up to 700 km, the
  source connects to a regional trunkline operating at an efficient
  scale (e.g., 3 Mtpa), paying only a proportional mass share of the
  CAPEX.
\item
  \textbf{The Trunk Leg (Booster Phase):} For distances exceeding 700
  km, the marginal cost of the trunkline increases by a penalty factor
  (\(P_{booster} = 2.0\)) to approximate the compounding capital and
  operational costs of intermediate booster stations.
\end{enumerate}

The trunkline CAPEX share (\(CAPEX_{trunk}\)) is calculated as:

\[CAPEX_{trunk} =
\begin{cases}
C_{base} \cdot (d_{effective} - 50) & \text{if } d_{effective} \le 700 \\
C_{base} \cdot (700 - 50) + (C_{base} \cdot P_{booster}) \cdot (d_{effective} - 700) & \text{if } d_{effective} > 700
\end{cases}\]

where \(C_{base}\) is the unit length cost of the efficient-scale
trunkline, scaled to the site's proportional mass.

\paragraph{Modality Optimization}\label{modality-optimization}

To eliminate artifactual spatial boundaries, the algorithm evaluates
both modalities for all onshore sinks and selects the optimal route:

\[Cost_{transport} = \min(C_{pipeline}, C_{shipping})\]

Because of the piecewise booster penalty, the pipeline cost curve
steepens significantly beyond 700 km. This ensures the pipeline cost
curve intersects with the flat \(C_{shipping}\) cost floor, allowing the
model to smoothly transition to shipping modalities at very long
distances.

\paragraph{Biochar Permanence and Agronomic Valuation (BEBCS
Feedbacks)}\label{biochar-permanence-and-agronomic-valuation-bebcs-feedbacks}

The Bioenergy with Biochar Storage (BEBCS) pathway diverts a fraction of
the original biomass carbon into a recalcitrant solid matrix. While
physical disposal methods such as the landfilling of biochar represent a
highly effective method to provide highly stable carbon storage, this
framework strictly evaluates soil application to quantify the cascading
agronomic co-benefits and regional substitution values. Consequently,
the carbon retention within the soil matrix is finite and must be
modeled dynamically based on localized biogeochemical conditions.

\subparagraph{\texorpdfstring{Biochar Carbon Permanence
(\(F_{perm}\))}{Biochar Carbon Permanence (F\_\{perm\})}}\label{biochar-carbon-permanence-f_perm}

To accurately quantify the net carbon dioxide removal (CDR) of the BEBCS
pathway, the model calculates the permanent carbon fraction
(\(F_{perm}\)) that remains stabilized in the soil after a 100-year time
horizon. This fraction is derived using the multi-pool decay methodology
established by (\citeproc{ref-woolf2021greenhouse}{\emph{11}}).

The baseline degradation rates of the biochar's labile and recalcitrant
carbon pools are fitted to either the peak pyrolysis temperature
(\(T_{py}\)) or the molar Hydrogen-to-Organic-Carbon (H:C) ratio of the
final char. Because microbial and chemical degradation kinetics are
highly temperature-dependent, these reference decay rates are spatially
adjusted to match the localized mean annual soil temperature using a
\(Q_{10}\) temperature sensitivity function. To optimize spatial
computational efficiency across the high-resolution rasters, the
thermodynamic module utilizes bilinear interpolation over a
pre-calculated matrix of \(F_{perm}\) values bounded by soil
temperatures (\(-55^\circ\text{C}\) to \(40^\circ\text{C}\)) and H:C
ratios (\(0.0\) to \(0.7\)).

Total carbon sequestered per megagram of feedstock is therefore modeled
as:

\[C_{seq} = Yield_{bc} \cdot C_{content} \cdot F_{perm} \cdot \left(\frac{44}{12}\right)\]

where \(Yield_{bc}\) is the mass fraction of biochar produced,
\(C_{content}\) is the carbon mass fraction of the biochar, and the
\(44/12\) ratio converts the elemental carbon to CO\textsubscript{2}
mass equivalents.

\paragraph{Mechanistic Agronomic
Valuation}\label{mechanistic-agronomic-valuation}

Previous spatial models frequently value biochar using a static,
exogenous global market price. To better reflect real-world agricultural
economics, C-SCAPE endogenously determines the economic credit of
biochar through a mechanistic substitution valuation. This assumes that
the rational economic value of biochar is equal to the cost of the
conventional agricultural inputs it successfully offsets:

\[Value_{ag} = Value_{lime} + Value_{nutrients} + Value_{physical}\]

\begin{enumerate}
\def\labelenumi{\arabic{enumi}.}
\tightlist
\item
  \textbf{Liming Equivalence (}\(Value_{lime}\)): The model spatially
  evaluates the native soil pH against a neutral agricultural target (pH
  6.5). If the soil is acidic, the biochar's Calcium Carbonate
  Equivalent (CCE, default 15\%) is credited against the localized
  market price of bulk agricultural lime (e.g., \(\$60/\text{Mg}\)). If
  the regional soil is already alkaline, the liming value is strictly
  zero.
\item
  \textbf{Nutrient Substitution (}\(Value_{nutrients}\)): The available
  nitrogen (N), phosphorus (P), and potassium (K) fractions within the
  biochar are credited against prevailing regional fertilizer costs.
  Baseline valuations rely on standard commercial substitutes, deriving
  elemental prices from Urea (\(\$0.92/\text{kg N}\)), Diammonium
  Phosphate (\(\$1.10/\text{kg P}_2\text{O}_5\)), and Potash
  (\(\$0.62/\text{kg K}_2\text{O}\)).
\item
  \textbf{Physical/CEC Improvements (}\(Value_{physical}\)): Biochar
  provides long-term yield efficiency gains by increasing water
  retention and buffering nutrient leaching, effects that are most
  pronounced in degraded or sandy soils. This value is heuristically
  scaled inversely to the soil's native Cation Exchange Capacity (CEC).
  Soils with severe CEC deficits receive higher annual premiums (up to
  \(\$50/\text{Mg}\)), while heavy clay soils (CEC
  \(> 30\text{ cmol/kg}\)) receive zero premium.
\end{enumerate}

Because physical and chemical soil improvements degrade or saturate over
time, the total annual physical and nutrient benefits are amortized as a
Present Value of an Annuity over a conservative 10-year agronomic impact
duration, discounted at the baseline financial discount rate.

\paragraph{Consequential Life Cycle Assessment and Climate
Mitigation}\label{consequential-life-cycle-assessment-and-climate-mitigation}

To robustly account for the climate mitigation benefit of bioenergy
generation, the framework evaluates greenhouse gas offsets using a
consequential Life Cycle Assessment (LCA) boundary. Total carbon
abatement is calculated distinctively for each technology pathway,
balancing active sequestration, grid displacement credits, and Scope 3
supply chain emissions.

\subparagraph{Net Carbon Abatement
Balances}\label{net-carbon-abatement-balances}

The net carbon abatement for each grid cell is defined as the sum of
physical carbon sequestered and emissions displaced, minus the
transportation emissions penalty:

\begin{itemize}
\tightlist
\item
  \textbf{BES:} Abatement relies entirely on the displacement of grid
  emissions, minus the Scope 3 emissions associated with diesel trucking
  of the biomass feedstock.
\item
  \textbf{BECCS:} Abatement includes the physical CO\textsubscript{2}
  captured and injected into the geological sink (assuming a 90\%
  capture rate) plus the grid displacement credit. However, because the
  parasitic energy load of the capture process reduces the net
  electrical efficiency (from 30\% to 28\%), the total energy
  exported---and thus the resulting grid displacement credit---is
  proportionally lower for BECCS than for pure BES{[}\^{}5{]}.
\item
  \textbf{BEBCS:} Abatement combines the grid displacement credit with
  the permanent CO\textsubscript{2} equivalent of the biochar carbon
  retained in the soil (\(F_{perm}\)), alongside a minor fixed soil
  greenhouse gas abatement credit (e.g., suppressed \(N_2O\) emissions).
\end{itemize}

For all pathways, Scope 3 feedstock transport emissions are modeled
dynamically by applying a heavy-duty diesel transport emissions factor
(\(0.0001\text{ Mg CO}_2\text{e}\) per \(\text{Mg}\cdot\text{km}\)) to
the effective biomass collection distance.

\paragraph{Grid Displacement Methodology (Marginal Carbon
Intensity)}\label{grid-displacement-methodology-marginal-carbon-intensity}

We estimate the Marginal Carbon Intensity (MCI) of displaced energy
under two bounding scenarios to reflect differing macro-economic climate
trajectories.

\paragraph{Scenario A: Empirical Baseline (Short-Run Marginal
Capacity)}\label{scenario-a-empirical-baseline-short-run-marginal-capacity}

Under a Business-as-Usual (BAU) trajectory, the model assumes that new
bioenergy capacity displaces the technology mix currently expanding on
regional grids. Using multi-country historical electricity supply data
(e.g., Ember, EIA, Eurostat), the model isolates generation sources
experiencing positive net growth (\(\Delta G > 0\)) over a trailing
five-year window (2018--2023).

The regional MCI is calculated as a generation-weighted average of these
expanding sources:

\[MCI = \frac{\sum(\Delta G_{i} \times CI_{i})}{\sum\Delta G_{i}}\]

where \(\Delta G_{i}\) represents the net positive generation added for
source \(i\), and \(CI_{i}\) represents the static Intergovernmental
Panel on Climate Change (IPCC) lifecycle greenhouse gas intensity for
that specific technology. Standard IPCC reference intensities utilized
include coal (820 \(\text{gCO}_2\text{eq/kWh}\)), natural gas (490
\(\text{gCO}_2\text{eq/kWh}\)), dedicated biomass (230
\(\text{gCO}_2\text{eq/kWh}\)), utility solar (48
\(\text{gCO}_2\text{eq/kWh}\)), wind (11.5
\(\text{gCO}_2\text{eq/kWh}\)), and nuclear (12
\(\text{gCO}_2\text{eq/kWh}\)). Applying this short-run marginal logic
yields empirical baseline MCIs of approximately 278
\(\text{gCO}_2\text{eq/kWh}\) for the United States and 528
\(\text{gCO}_2\text{eq/kWh}\) for China.

\paragraph{Scenario B: Paris-Aligned Counterfactual (Long-Run Structural
Shift)}\label{scenario-b-paris-aligned-counterfactual-long-run-structural-shift}

To account for structural shifts in energy systems operating under
strict climate stabilization limits (e.g., 1.5°C or 2.0°C targets), a
secondary counterfactual scenario is evaluated. In a deep mitigation
regime, fossil fuels are systematically phased out by hard policy caps.
Therefore, the deployment of a baseload negative-emissions technology
like BECCS or BEBCS does not offset high-emission fossil generation;
rather, it displaces or delays the deployment of the next most
capital-intensive low-carbon baseload alternative.

Under this sensitivity mode, the model assumes the displaced marginal
generation is nuclear power. This scenario aggressively compresses net
mitigation margins, dropping the grid displacement credit to a
standardized lifecycle factor of 12 \(\text{gCO}_2\text{eq/kWh}\)
globally.

\subsubsection{Model Parameters and Data
Sources}\label{model-parameters-and-data-sources}

To ensure full reproducibility and transparency, the operational,
financial, and biophysical parameters utilized within the C-SCAPE
framework are tabulated below. Parameters are classified by their
spatial resolution: \textbf{Scalar} values are held constant globally;
\textbf{Regional} values act as constants but are overridden based on
the macro-economic domain (e.g., India, China, Europe, US); and
\textbf{Spatial} values are resolved dynamically at the grid-cell level
using high-resolution raster datasets.

\begingroup
\fontsize{7.0pt}{8.0pt}\selectfont

\begin{longtable}{>{\raggedright\arraybackslash}p{\dimexpr 0.18\linewidth -2\tabcolsep-1.5\arrayrulewidth}>{\raggedright\arraybackslash}p{\dimexpr 0.24\linewidth -2\tabcolsep-1.5\arrayrulewidth}>{\raggedright\arraybackslash}p{\dimexpr 0.12\linewidth -2\tabcolsep-1.5\arrayrulewidth}>{\raggedleft\arraybackslash}p{\dimexpr 0.10\linewidth -2\tabcolsep-1.5\arrayrulewidth}>{\raggedleft\arraybackslash}p{\dimexpr 0.10\linewidth -2\tabcolsep-1.5\arrayrulewidth}>{\raggedleft\arraybackslash}p{\dimexpr 0.10\linewidth -2\tabcolsep-1.5\arrayrulewidth}>{\raggedright\arraybackslash}p{\dimexpr 0.16\linewidth -2\tabcolsep-1.5\arrayrulewidth}}

\caption{\label{tbl-parameters}Operational, financial, and biophysical
parameters utilized within the BiocharAG framework.}

\tabularnewline

\toprule
 &  & \multicolumn{4}{>{\centering\arraybackslash}m{\dimexpr 0.42\linewidth -2\tabcolsep-1.5\arrayrulewidth}}{\parbox{\linewidth}{\centering {Value}}} &  \\ 
\cmidrule(lr){3-6}
Parameter & Description & USA/Default & Europe & China & India & Units \\ 
\midrule\endhead\addlinespace[2.5pt]
\multicolumn{7}{l}{Financial \& Market} \\[2.5pt] 
\midrule\addlinespace[2.5pt]
\texttt{discount\_rate} & Cost of capital & 0.05 & 0.045 & 0.045 & 0.1 & Fraction \\ 
\texttt{elec\_price} & Wholesale electricity price & Spatially explicit &  &  &  & USD / MWh \\ 
\texttt{wholesale\_discount} & Ratio of wholesale to retail price & 0.4 & 0.4 & 0.4 & 1 & Fraction \\ 
\texttt{c\_price} & Carbon abatement credit price & 50 &  &  &  & USD / tCO\textsubscript{2}e \\ 
\texttt{bes\_life} / \texttt{beccs\_life} & Infrastructure project lifetime & 30 &  &  &  & Years \\ 
\texttt{capacity\_factor} & Annual plant capacity utilization & 0.85 &  &  &  & Fraction \\ 
\texttt{scaling\_factor} & Capital cost scale exponent & 0.7 &  &  &  & --- \\ 
\midrule\addlinespace[2.5pt]
\multicolumn{7}{l}{Technology \& Thermodynamics} \\[2.5pt] 
\midrule\addlinespace[2.5pt]
\texttt{plant\_mw\_th} & Target plant thermal capacity & 250 &  &  &  & MWth \\ 
\texttt{bes\_capital\_cost} & Base CAPEX for BES & 3750 & 3450 & 2100 & 1950 & USD / kWe \\ 
\texttt{beccs\_capital\_cost} & Base CAPEX for BECCS & 5000 & 4600 & 2800 & 2600 & USD / kWe \\ 
\texttt{bes\_om\_factor} & O\&M cost as fraction of CAPEX & 0.045 & 0.045 & 0.03 & 0.025 & Fraction \\ 
\texttt{beccs\_om\_factor} & O\&M cost for BECCS & 0.055 & 0.055 & 0.035 & 0.03 & Fraction \\ 
\texttt{bes\_energy\_efficiency} & Electrical efficiency for BES & 0.3 &  &  &  & Fraction \\ 
\texttt{beccs\_efficiency} & Electrical efficiency for BECCS & 0.28 &  &  &  & Fraction \\ 
\texttt{capture\_rate} & BECCS CO2 capture efficiency & 0.9 &  &  &  & Fraction \\ 
\texttt{py\_temp} & Peak pyrolysis temperature & 500 &  &  &  & °C \\ 
\texttt{heat\_loss} & Pyrolysis process heat loss & 2.35 &  &  &  & GJ / Mg \\ 
\texttt{bo\_hhv} & Bio-oil higher heating value & 27.63 &  &  &  & GJ / Mg \\ 
\texttt{bo\_c}, \texttt{bo\_h}, \texttt{bo\_o} & Bio-oil elemental mass fractions & 0.625, 0.0756, 0.2994 &  &  &  & Fraction \\ 
\midrule\addlinespace[2.5pt]
\multicolumn{7}{l}{Feedstock \& Logistics} \\[2.5pt] 
\midrule\addlinespace[2.5pt]
\texttt{bm\_lhv} & Biomass lower heating value & 18.6 &  &  &  & GJ / Mg \\ 
\texttt{bm\_c} & Biomass carbon mass fraction & 0.48 &  &  &  & Fraction \\ 
\texttt{bm\_ash} & Biomass ash mass fraction & 0.05 & 0.05 & 0.05 & 0.15 & Fraction \\ 
\texttt{bm\_transport\_fixed} & Fixed loading/handling cost & 5 &  &  &  & USD / Mg \\ 
\texttt{bm\_transport\_var} & Variable trucking cost & 0.15 &  &  &  & USD / Mg / km \\ 
\texttt{transport\_EF} & Heavy-duty truck emissions & 0.0001 &  &  &  & Mg CO\textsubscript{2}e / Mg·km \\ 
\midrule\addlinespace[2.5pt]
\multicolumn{7}{l}{CO\textasciitilde{}2\textasciitilde{} Transport Infrastructure} \\[2.5pt] 
\midrule\addlinespace[2.5pt]
\texttt{ccs\_storage\_cost} & Onshore injection and monitoring & 12 &  &  &  & USD / tCO\textsubscript{2} \\ 
\texttt{cost\_offshore\_storage} & Offshore injection and monitoring & 40 &  &  &  & USD / tCO\textsubscript{2} \\ 
\texttt{feeder\_threshold} & Dedicated feeder pipeline limit & 50 &  &  &  & km \\ 
\texttt{booster\_threshold} & Distance triggering booster penalty & 700 &  &  &  & km \\ 
\texttt{booster\_penalty} & Marginal cost multiplier for boosters & 2 &  &  &  & --- \\ 
\texttt{trunk\_mass\_flow} & Ref. capacity for trunkline sharing & 3000000 &  &  &  & Mg / year \\ 
\texttt{liq\_cost} & CO\textsubscript{2} Liquefaction cost (shipping) & 20 &  &  &  & USD / tCO\textsubscript{2} \\ 
\texttt{term\_cost} & Terminal handling cost (shipping) & 15 &  &  &  & USD / tCO\textsubscript{2} \\ 
\texttt{voyage\_rate} & Variable voyage cost (shipping) & 0.035 &  &  &  & USD / tCO\textsubscript{2} / km \\ 
\texttt{transport\_elec\_price} & Electricity price for compression & 0.1 &  &  &  & USD / kWh \\ 
\midrule\addlinespace[2.5pt]
\multicolumn{7}{l}{Biochar \& Agronomic} \\[2.5pt] 
\midrule\addlinespace[2.5pt]
\texttt{h\_c\_org} & Molar H:C ratio for permanence & 0.35 &  &  &  & Molar Ratio \\ 
\texttt{time\_frame} & Time horizon of pyC permanence & 100 &  &  &  & Years \\ 
\texttt{target\_ph} & Target agricultural soil pH & 6.5 &  &  &  & pH \\ 
\texttt{bc\_cce} & Calcium Carbonate Equivalent & 0.15 &  &  &  & Fraction \\ 
\texttt{price\_lime} & Market price for agricultural lime & 45 & 50 & 35 & 35 & USD / Mg \\ 
\texttt{price\_n} & Market price for Nitrogen & 1.59 & 1.75 & 0.79 & 0.14 & USD / kg N \\ 
\texttt{price\_p} & Market price for Phosphorus & 2.08 & 2.29 & 1.1 & 0.7 & USD / kg P\textsubscript{2}O\textsubscript{5} \\ 
\texttt{price\_k} & Market price for Potassium & 0.82 & 0.9 & 0.55 & 0.68 & USD / kg K\textsubscript{2}O \\ 
\texttt{avail-n / -p / -k} & First-year nutrient availability & 0.1, 0.5, 0.8 &  &  &  & Fraction \\ 
\texttt{ag\_impact\_duration} & Duration of nutrient/liming benefit & 5 &  &  &  & Years \\ 
\texttt{soil\_ghg\_abatement} & Suppressed soil emissions credit & 0.1 &  &  &  & tCO\textsubscript{2}e / Mg \\ 
\midrule\addlinespace[2.5pt]
\multicolumn{7}{l}{Spatial \& Environmental} \\[2.5pt] 
\midrule\addlinespace[2.5pt]
\texttt{ff\_c\_intensity} & Displaced grid carbon intensity & Spatially explicit &  &  &  & tCO\textsubscript{2}eq / GJ \\ 
\texttt{biomass\_density} & Available spatial crop residues & Spatially explicit &  &  &  & Mg / km\textsuperscript{2} \\ 
\texttt{soil\_ph} & Native topsoil pH & Spatially explicit &  &  &  & pH \\ 
\texttt{soil\_cec} & Cation Exchange Capacity & Spatially explicit &  &  &  & cmol / kg \\ 
\texttt{soil\_temp} & Mean annual soil temperature & Spatially explicit &  &  &  & °C \\ 
\bottomrule

\end{longtable}

\endgroup

\section{Results}\label{results}

\subsection{Capital Intensive Pathways are Highly Sensitive to Discount
Rates and Carbon
Prices}\label{capital-intensive-pathways-are-highly-sensitive-to-discount-rates-and-carbon-prices}

Figures \hyperref[fig-evap-usa]{fig-evap-usa} to
\hyperref[fig-evap-china]{fig-evap-china} illustrate the core
vulnerability of CDR pathways to capital lock-out. Because both BECCS
(heavy upfront infrastructure) and BEBCS (deferred agronomic value) are
highly sensitive to the cost of capital, high discount rates act as a
strong barrier to adoption, forcing a retreat to traditional Bioenergy
(BES) without high carbon prices. For example, in the USA, at a 2\%
social discount rate and a \$30/t carbon price, BEBCS covers a large
extent of the US Midwest and South, with BECCS only being viable close
to CO\textsubscript{2} sinks. However, stepping up to the 8\% market
baseline eliminates BEBCS's advantage, replacing it with almost entirely
with BES. The value of biochar relies heavily on liming and nutrient
substitution benefits that are amortized over a 10-year agronomic impact
duration. At an 8\% or 15\% discount rate, the present value of those
future soil benefits is diminished, rendering the upfront thermal
penalties of pyrolysis economically unviable at low carbon prices. At a
2\% discount rate and intermediate carbon price (around \$100/t), BECCS
is nearly universal across the US. At 15\%, the BECCS dominance is
fractured, with BES reclaiming the geographic margins (e.g., the
Northeast and far West), and BEBCS remaining only in localized pockets
where soil conditions make it particularly attractive. This occurs
because, when capital is expensive, the penalty for long-distance
pipelines becomes insurmountable until carbon prices reach the \$150/t
range. By the time the carbon price reaches \$150/t, the maps become
almost entirely red across all discount rates. At this pricing level,
the high CDR capture efficiency overwhelms the high cost of capital,
making BECCS the most economically attractive option.

\subsection{Biochar provides a mid-price transition technology option
and a long term solution for specific agro-ecological
zones}\label{biochar-provides-a-mid-price-transition-technology-option-and-a-long-term-solution-for-specific-agro-ecological-zones}

Figure \hyperref[fig-evap-china]{fig-evap-china} provides a counterpoint
to the US results. It shows how soils with strong production constraints
that can be alleviated by biochar can create a persistent competitive
advantage for biochar, particularly when combined with large distances
to CO\textsubscript{2} storage sites. BEBCS performs a central role as a
transition technology that peaks at intermediate carbon prices that are
high enough to make its carbon sequestration competitive with BES, but
not so high that the more efficient CDR of BECCS per unit biomass
outcompetes it (Figures \hyperref[fig-MACC-eur]{fig-MACC-eur},
\hyperref[fig-MACC-india]{fig-MACC-india}, and
\hyperref[fig-MACC-china]{fig-MACC-china}). Nonetheless, while BECCS
ultimately becomes the optimal technology across the entire US spatial
domain at high carbon prices, BEBCS remains the most viable pathway in
southern China across almost all price and discount rate scenarios. This
demonstrates that BEBCS is not merely a transition technology, but the
optimal long-term pathway for specific agro-ecological zones. In
Southern China, characterized by highly weathered, acidic soils, the
biochar's liming and nutrient use efficiency benefits provide a
substantial economic credit that buffers against high processing costs.
The model scales physical yield efficiency gains inversely to the soil's
native Cation Exchange Capacity (CEC), thus granting severe CEC deficits
higher premiums. This is confirmed by the SHAP analysis (Figure
\hyperref[fig-shap-dominant]{fig-shap-dominant}), which identifies soil
CEC as the dominant spatial driver in this region.

\subsection{Regional Policies Hinder Biochar
Competitiveness}\label{regional-policies-hinder-biochar-competitiveness}

In India, on the other hand, the mid carbon price advantage of BEBCS is
almost entirely eliminated, with BECCS eventually covering most of the
map at high carbon prices. Here, fertilizer subsidies hamper BEBCS's
economic viability. Because BEBCS economics relies on improved nutrient
use efficiency and lower fertilizer inputs per unit crop yield, cheap
subsidized synthetic fertilizers erode biochar's nutrient substitution
value. Despite this, Biochar maintains a strong foothold in the regions
in the North East, North West, and central Eastern coastal regions of
India where low soil CEC is a major constraint to crop production.

\subsection{Regional Energy prices and Carbon Intensity provide a
critical constraint on the optimal CDR
pathways}\label{regional-energy-prices-and-carbon-intensity-provide-a-critical-constraint-on-the-optimal-cdr-pathways}

At lower carbon prices The optimal technology in Europe is almost
universally BES, driven by high energy prices as the dominant driver
across Europe. At \$30 / t CO\textsubscript{2}, the revenue from
maximized power generation in BES outweighs the incentive to capture and
sequester carbon, rendering the capital expenditure for BECCS and BEBCS
unviable across all discount rates. As carbon prices increase, BECCS
begins to dominate across central and northern Europe. BEBCS gains a
strong advantage in Southern Europe, and Eastern Scandinavia where
various combinations of large distances to CO\textsubscript{2} sinks,
already-low carbon intensity of the energy systems, and infertile soils
make it more competitive. As in China, we see BEBCS maintain its
advantage in these regions, even at very high carbon prices.

\section{Conclusion}\label{conclusion}

The deployment of biomass for climate change mitigation cannot be guided
by a single universal pathway. Our spatial analysis demonstrates that
the economic viability of BES, BECCS, and BEBCS is dictated by a complex
interplay of localized biophysical constraints and regional market
dynamics.

While BECCS offers the highest theoretical carbon removal potential, its
heavy infrastructure requirements and sensitivity to the cost of capital
create significant barriers to entry, particularly where biomass is
located far from suitable geological sinks. In contrast, biochar (BEBCS)
circumvents these transport constraints and establishes a resilient
economic niche in regions with degraded, low-CEC soils, where its
agronomic co-benefits---such as liming equivalence and nutrient
substitution---offset its lower electrical conversion efficiencies.
Southern China exemplifies this dynamic, demonstrating that BEBCS is a
robust, long-term mitigation strategy for specific agro-ecological zones
rather than merely a transition technology.

Furthermore, baseline energy markets and agricultural policies heavily
distort optimal pathways. In Europe, high wholesale electricity prices
establish BES as the dominant pathway at low carbon prices, while in
India, fertilizer subsidies artificially suppress biochar's agronomic
value, pushing the market toward traditional BES.

Realizing the climate mitigation potential of global biomass requires
moving beyond blanket technological mandates. Policymakers must adopt
spatially explicit strategies that target BECCS deployment near
geological sinks, incentivize BEBCS in agro-ecological zones suffering
from soil degradation, and recognize the enduring role of BES in
displacing highly carbon-intensive fossil grids.

\section{Figures}\label{figures}

\begin{figure}

\centering{

\includegraphics[width=1\linewidth,height=\textheight,keepaspectratio]{images/Global_Fig8_Breakeven_CPrice_CP100_MW250_reg.png}

}

\caption{\label{fig-break-even}Minimum Break-even carbon price for each
technology option. The technology with the lowest break-even price is
shown as the ``best technology'' in the penultimate row. The final row
shows the break-even price of this technology.}

\end{figure}%

\begin{figure}

\centering{

\includegraphics[width=1\linewidth,height=\textheight,keepaspectratio]{images/US_Fig3_Evaporation_Maps_CP100_MW250.png}

}

\caption{\label{fig-evap-usa}Optimal technology as a function of carbon
price and discount rate (USA).}

\end{figure}%

\begin{figure}

\centering{

\includegraphics[width=1\linewidth,height=\textheight,keepaspectratio]{images/Europe_Fig3_Evaporation_Maps_CP100_MW250.png}

}

\caption{\label{fig-evap-eur}Optimal technology as a function of carbon
price and discount rate (Europe).}

\end{figure}%

\begin{figure}

\centering{

\includegraphics[width=1\linewidth,height=\textheight,keepaspectratio]{images/US_Fig3_Evaporation_Maps_CP100_MW250_reg.png}

}

\caption{\label{fig-evap-india}Optimal technology as a function of
carbon price and discount rate (India).}

\end{figure}%

\begin{figure}

\centering{

\includegraphics[width=1\linewidth,height=\textheight,keepaspectratio]{images/China_Fig3_Evaporation_Maps_CP100_MW250_reg.png}

}

\caption{\label{fig-evap-china}Optimal technology as a function of
carbon price and discount rate (China).}

\end{figure}%

\begin{figure}

\centering{

\includegraphics[width=0.5\linewidth,height=\textheight,keepaspectratio]{images/Europe_Fig6_MACC_CP100_MW250.png}

}

\caption{\label{fig-MACC-eur}Marginal abatement cost curve for Europe.}

\end{figure}%

\begin{figure}

\centering{

\includegraphics[width=0.5\linewidth,height=\textheight,keepaspectratio]{images/India_Fig6_MACC_EA_CP100_MW250.png}

}

\caption{\label{fig-MACC-india}Marginal abatement cost curve for India.}

\end{figure}%

\begin{figure}

\centering{

\includegraphics[width=0.5\linewidth,height=\textheight,keepaspectratio]{images/China_Fig6_MACC_EA_CP100_MW250.png}

}

\caption{\label{fig-MACC-china}Marginal abatement cost curve for China.}

\end{figure}%

\begin{figure}

\centering{

\includegraphics[width=1\linewidth,height=\textheight,keepaspectratio]{images/map_dominant_shap_feature.png}

}

\caption{\label{fig-shap-dominant}Dominant driver of breakeven carbon
prices (Shapley decomposition).}

\end{figure}%

\begin{figure}

\centering{

\includegraphics[width=1\linewidth,height=\textheight,keepaspectratio]{images/map_max_shap_magnitude-2.png}

}

\caption{\label{fig-shap-magnitude}Magnitude of the dominant driver
(Shapley decomposition).}

\end{figure}%

\section{References}\label{references}

\phantomsection\label{refs}
\begin{CSLReferences}{0}{1}
\bibitem[\citeproctext]{ref-woolf2010}
\CSLLeftMargin{1. }%
\CSLRightInline{D. Woolf, J. E. Amonette, F. A. Street-Perrott, J.
Lehmann, S. Joseph,
\href{https://doi.org/10.1038/ncomms1053}{Sustainable biochar to
mitigate global climate change}. \emph{Nature Communications}
\textbf{1}, 56 (2010).}

\bibitem[\citeproctext]{ref-lehmann2021biochar}
\CSLLeftMargin{2. }%
\CSLRightInline{J. Lehmann, A. Cowie, A. Allison, M. Brandão, S. Joseph,
D. Powlson, D. Woolf,
\href{https://doi.org/10.1038/s41561-021-00852-8}{Biochar in climate
change mitigation}. \emph{Nature Geoscience} \textbf{14}, 883--892
(2021).}

\bibitem[\citeproctext]{ref-woolf2016}
\CSLLeftMargin{3. }%
\CSLRightInline{D. Woolf, J. Lehmann, D. R. Lee,
\href{https://doi.org/10/f87kfp}{Optimal bioenergy power generation for
climate change mitigation with or without carbon sequestration}.
\emph{Nature Communications} \textbf{7}, 13160 (2016).}

\bibitem[\citeproctext]{ref-peters2015}
\CSLLeftMargin{4. }%
\CSLRightInline{J. F. Peters, D. Iribarren, J. Dufour,
\href{https://doi.org/10.1021/es5060786}{Biomass pyrolysis for biochar
or energy applications? A life cycle assessment}. \emph{Environmental
Science \& Technology} \textbf{49}, 5195--5202 (2015).}

\bibitem[\citeproctext]{ref-lefebvre2021}
\CSLLeftMargin{5. }%
\CSLRightInline{D. Lefebvre, A. Williams, G. J. D. Kirk, J. Meersmans,
S. Sohi, P. Goglio, P. Smith,
\href{https://doi.org/10.1016/j.jclepro.2021.127764}{An anticipatory
life cycle assessment of the use of biochar from sugarcane residues as a
greenhouse gas removal technology}. \emph{Journal of Cleaner Production}
\textbf{312}, 127764 (2021).}

\bibitem[\citeproctext]{ref-karan2023}
\CSLLeftMargin{6. }%
\CSLRightInline{S. K. Karan, D. Woolf, E. S. Azzi, C. Sundberg, S. A.
Wood, \href{https://doi.org/10.1111/gcbb.13102}{Potential for biochar
carbon sequestration from crop residues: A global spatially explicit
assessment}. \emph{GCB Bioenergy} \textbf{15}, 1424--1436 (2023).}

\bibitem[\citeproctext]{ref-sultana2014}
\CSLLeftMargin{7. }%
\CSLRightInline{A. Sultana, A. Kumar,
\href{https://doi.org/10.1016/j.apenergy.2013.12.036}{Development of
tortuosity factor for assessment of lignocellulosic biomass delivery
cost to a biorefinery}. \emph{Applied Energy} \textbf{119}, 288--295
(2014).}

\bibitem[\citeproctext]{ref-luan2025a}
\CSLLeftMargin{8. }%
\CSLRightInline{J. Luan, Q. Wang, D. Shao, B. Cui, P. Han, Q. He,
\href{https://doi.org/10.3389/fenrg.2025.1634354}{Research progress on
influencing factors and control methods of slagging in biomass
combustion}. \emph{Frontiers in Energy Research} \textbf{13} (2025).}

\bibitem[\citeproctext]{ref-baxter1998}
\CSLLeftMargin{9. }%
\CSLRightInline{L. L. Baxter, T. R. Miles, T. R. Miles, B. M. Jenkins,
T. Milne, D. Dayton, R. W. Bryers, L. L. Oden,
\href{https://doi.org/10.1016/S0378-3820(97)00060-X}{The behavior of
inorganic material in biomass-fired power boilers: Field and laboratory
experiences}. \emph{Fuel Processing Technology} \textbf{54}, 47--78
(1998).}

\bibitem[\citeproctext]{ref-garcuxeda2017}
\CSLLeftMargin{10. }%
\CSLRightInline{R. García, C. Pizarro, A. G. Lavín, J. L. Bueno,
\href{https://doi.org/10.1016/j.fuel.2017.01.063}{Biomass sources for
thermal conversion. Techno-economical overview}. \emph{Fuel}
\textbf{195}, 182--189 (2017).}

\bibitem[\citeproctext]{ref-woolf2021greenhouse}
\CSLLeftMargin{11. }%
\CSLRightInline{D. Woolf, J. Lehmann, S. Ogle, A. W. Kishimoto-Mo, B.
McConkey, J. Baldock,
\href{https://doi.org/10.3389/fclim.2021.684182}{Greenhouse gas removal
and emissions negative technologies}. \emph{Frontiers in Climate}
\textbf{3}, 684182 (2021).}

\end{CSLReferences}




\end{document}
